% ============================================================
%  Sprawozdanie: Badanie dyfrakcji światła (DYFRA)
%  Fotonika – laboratorium, ZUT Szczecin
% ============================================================
\documentclass[a4paper,12pt]{article}

\usepackage[utf8]{inputenc}
\usepackage[T1]{fontenc}
\usepackage{lmodern} % <-- DODANO: Wektorowa, gładka czcionka z ładnymi polskimi znakami
\usepackage[polish]{babel}
\usepackage{geometry}
\geometry{a4paper, top=2.2cm, bottom=2.2cm, left=2.5cm, right=2.5cm}

\usepackage{amsmath,amssymb}
\usepackage{booktabs}
\usepackage{graphicx}
\usepackage{float}
\usepackage{caption}
\usepackage{multirow}
\usepackage{array}
\usepackage{xcolor}
\usepackage{hyperref}
\usepackage{siunitx}
\sisetup{output-decimal-marker={,}, locale=DE}
\usepackage{parskip}
\setlength{\parskip}{4pt}

\title{}
\date{}

\begin{document}

% ============================================================
%  TABELA NAGŁÓWKOWA (punkty 1–3)
% ============================================================
\begin{table}[H]
\centering
\small
\renewcommand{\arraystretch}{1.3}
\resizebox{\textwidth}{!}{% <-- Skalowanie tabeli
\begin{tabular}{|*{12}{c|}}
\hline
\multicolumn{12}{|c|}{\textbf{Zachodniopomorski Uniwersytet Technologiczny w Szczecinie}} \\
\hline
\multicolumn{2}{|c|}{Przedmiot:} & \multicolumn{10}{c|}{Fotonika} \\
\hline
\multicolumn{2}{|c|}{Prowadzący:} & \multicolumn{10}{c|}{Mgr. Inż. Eliza Miśkiewicz} \\
\hline
\multicolumn{12}{|c|}{\textbf{SPRAWOZDANIE Z WYKONANEGO ĆWICZENIA}} \\
\hline
\multicolumn{2}{|c|}{Nr ćw.: 7} & 04 & \multicolumn{2}{c|}{Temat:} & \multicolumn{7}{c|}{Pomiar stężenia substancji w roztworze (LAMB)} \\
\hline
Zespół lab.: & 5 & Grupa: & L1 & Studia: & S2 & Kierunek: & TI & Rok: & 2025/2026 & Data: & 20.04.2026 \\
\hline
\multicolumn{3}{|c|}{Skład zespołu wg. obecności podczas ćwiczenia:} & \multicolumn{9}{c|}{1. Kamil Suchy}{2. Piotr Gadomski} \\
\hline
\multicolumn{2}{|c|}{Ocena:} & \multicolumn{10}{c|}{} \\
\hline
\end{tabular}%
}
\end{table}
